\begin{proof}[Proof of \cref{obj:lem:cai-low-chebyshev-approximation}]
\begin{enumerate}
\item Fix an even integer \(D\ge2\), and set
\[
\ell_{\mathrm{tr}}:=D/2\in\mathbb N,\qquad
\ell_{\mathrm{tr}}\ge1,\qquad D=2\ell_{\mathrm{tr}}.
\]
For the truncation argument, extend the polynomial family to degree zero by
\[
q_0(x):=\frac2\pi.
\]
Thus, for every \(\ell\in\mathbb N_0\),
\[
q_{2\ell}(x)
=\frac2\pi+\frac4\pi\sum_{\ell'=1}^{\ell}
 \frac{(-1)^{\ell'+1}\mathcal T_{2\ell'}(x)}
 {4(\ell')^2-1},
\qquad q_D=q_{2\ell_{\mathrm{tr}}}.
\]
This is precisely the one-based truncation of the even Chebyshev modes; the sum is empty when \(\ell=0\).
% lean: CausalSmith.Stat.LdpOptvalueUniformFrontier.chebyshevAbsPoly_eq_library

\item Define the continuous, \(\pi\)-periodic function \(\varphi_{\sin}:\mathbb R\to\mathbb R\) by
\[
\varphi_{\sin}(\vartheta)=\sin\vartheta
\quad(0\le\vartheta\le\pi),\qquad
\varphi_{\sin}(\vartheta+\pi)=\varphi_{\sin}(\vartheta)
\quad(\vartheta\in\mathbb R).
\]
The endpoint values agree, so the periodic extension is continuous. With \(\mathrm i:=\sqrt{-1}\), define its Fourier coefficients by
\[
\widehat\varphi_{\sin}(\ell)
:=\frac1\pi\int_0^\pi
 \sin\vartheta\,e^{-2\mathrm i\ell\vartheta}\,d\vartheta,
\qquad \ell\in\mathbb Z.
\]
For each such \(\ell\), \(1-4\ell^2\ne0\), and the function
\[
\Phi_\ell(\vartheta)
:=e^{-2\mathrm i\ell\vartheta}
 \frac{-2\mathrm i\ell\sin\vartheta-\cos\vartheta}
 {1-4\ell^2},
\qquad \vartheta\in\mathbb R,
\]
satisfies
\[
\Phi_\ell'(\vartheta)
=e^{-2\mathrm i\ell\vartheta}\sin\vartheta.
\]
Since \(e^{-2\mathrm i\ell\pi}=1\), evaluation at the endpoints gives
\[
\widehat\varphi_{\sin}(\ell)
=\frac{\Phi_\ell(\pi)-\Phi_\ell(0)}{\pi}
=\frac{2}{\pi(1-4\ell^2)}
=-\frac{2}{\pi(4\ell^2-1)}.
\]
In particular, \(\widehat\varphi_{\sin}(0)=2/\pi\).
% lean: Causalean.Mathlib.Analysis.AbsoluteValueMomentPriorDuality.fourierCoeff_sinCircle

\item For every integer \(\ell\ge1\),
\[
\bigl|\widehat\varphi_{\sin}(\ell)\bigr|
=\frac{2}{\pi(4\ell^2-1)}
\le\frac1{\ell^2}.
\]
Indeed, \(\ell^2\ge1\) and \(\pi\ge2\) imply
\[
2\ell^2\le\pi(4\ell^2-1).
\]
The coefficients at \(\ell\) and \(-\ell\) agree. Comparison with the convergent series \(\sum_{\ell=1}^\infty\ell^{-2}\) therefore yields
\[
\sum_{\ell\in\mathbb Z}
 \bigl|\widehat\varphi_{\sin}(\ell)\bigr|<\infty.
\]
The Fourier reconstruction theorem for a continuous periodic function with absolutely summable coefficients now gives
\[
\varphi_{\sin}(\vartheta)
=\sum_{\ell\in\mathbb Z}
 \widehat\varphi_{\sin}(\ell)e^{2\mathrm i\ell\vartheta},
\qquad \vartheta\in\mathbb R.
\]
% lean: Causalean.Mathlib.Analysis.AbsoluteValueMomentPriorDuality.summable_fourierCoeff_sinCircle

\item For \(0\le\vartheta\le\pi\), periodicity and the sine addition identities give
\[
\varphi_{\sin}(\vartheta+\pi/2)
=
\begin{cases}
\cos\vartheta,&0\le\vartheta\le\pi/2,\\
-\cos\vartheta,&\pi/2<\vartheta\le\pi,
\end{cases}
=|\cos\vartheta|.
\]
The defining Chebyshev recurrence and the cosine addition formula also give
\[
\mathcal T_\ell(\cos\vartheta)=\cos(\ell\vartheta),
\qquad \ell\in\mathbb N_0,\quad \vartheta\in\mathbb R.
\]
Consequently, for \(\ell\ge1\), the paired Fourier terms at the shifted argument satisfy
\[
\begin{aligned}
&\widehat\varphi_{\sin}(\ell)
 e^{2\mathrm i\ell(\vartheta+\pi/2)}
+\widehat\varphi_{\sin}(-\ell)
 e^{-2\mathrm i\ell(\vartheta+\pi/2)}
\\
&\qquad
=-\frac{4\cos(2\ell\vartheta+\ell\pi)}
 {\pi(4\ell^2-1)}
=\frac4\pi
 \frac{(-1)^{\ell+1}\mathcal T_{2\ell}(\cos\vartheta)}
 {4\ell^2-1}.
\end{aligned}
\]
For \(x\in[-1,1]\), define
\[
\vartheta_x:=\arccos x\in[0,\pi],
\qquad \cos\vartheta_x=x.
\]
Pairing the positive and negative frequencies in the absolutely convergent Fourier series and taking real parts therefore yields
\[
|x|=\frac2\pi+\frac4\pi\sum_{\ell=1}^{\infty}
 \frac{(-1)^{\ell+1}\mathcal T_{2\ell}(x)}
 {4\ell^2-1}.
\]
% lean: Causalean.Mathlib.Analysis.AbsoluteValueMomentPriorDuality.hasSum_absChebTerm_real

\item Splitting this series at \(\ell_{\mathrm{tr}}\) gives
\[
|x|-q_D(x)
=\frac4\pi\sum_{\ell=\ell_{\mathrm{tr}}+1}^{\infty}
 \frac{(-1)^{\ell+1}\mathcal T_{2\ell}(x)}
 {4\ell^2-1}.
\]
The scalar majorant has an exact telescoping sum, because
\[
\frac1{4\ell^2-1}
=\frac12\left(\frac1{2\ell-1}-\frac1{2\ell+1}\right)
\qquad(\ell\ge1).
\]
In particular, for every integer \(\ell_{\mathrm{end}}\ge\ell_{\mathrm{tr}}\),
\[
\sum_{\ell=\ell_{\mathrm{tr}}+1}^{\ell_{\mathrm{end}}}
 \frac1{4\ell^2-1}
=\frac12\left(
 \frac1{2\ell_{\mathrm{tr}}+1}
 -\frac1{2\ell_{\mathrm{end}}+1}\right),
\]
including the empty sum when \(\ell_{\mathrm{end}}=\ell_{\mathrm{tr}}\). Passing to the limit gives
\[
\sum_{\ell=\ell_{\mathrm{tr}}+1}^{\infty}
 \frac1{4\ell^2-1}
=\frac1{2(2\ell_{\mathrm{tr}}+1)}.
\]
Moreover, the cosine identity implies
\[
|\mathcal T_{2\ell}(x)|
=|\cos(2\ell\vartheta_x)|\le1.
\]
Thus the residual series is absolutely convergent, and
\[
\begin{aligned}
\bigl|q_D(x)-|x|\bigr|
&\le\frac4\pi
 \sum_{\ell=\ell_{\mathrm{tr}}+1}^{\infty}
 \frac{|\mathcal T_{2\ell}(x)|}{4\ell^2-1}\\
&\le\frac4\pi
 \sum_{\ell=\ell_{\mathrm{tr}}+1}^{\infty}
 \frac1{4\ell^2-1}
=\frac2{\pi(2\ell_{\mathrm{tr}}+1)}
=\frac2{\pi(D+1)}.
\end{aligned}
\]
% lean: Causalean.Mathlib.Analysis.AbsoluteValueMomentPriorDuality.absChebPoly_error_le

\item For real polynomials
\[
\mathcal B(x)=\sum_{\ell\ge0}\beta_\ell x^\ell,
\qquad
\mathcal E(x)=\sum_{\ell\ge0}\gamma_\ell x^\ell,
\]
where the real coefficient sequences have finite support, define
\[
\|\mathcal B\|_{\mathrm{coef},1}
:=\sum_{\ell\ge0}|\beta_\ell|,
\qquad
\|\mathcal E\|_{\mathrm{coef},1}
:=\sum_{\ell\ge0}|\gamma_\ell|.
\]
Coefficientwise triangle inequalities give
\[
\|\mathcal B\pm\mathcal E\|_{\mathrm{coef},1}
\le
\|\mathcal B\|_{\mathrm{coef},1}
+\|\mathcal E\|_{\mathrm{coef},1}.
\]
For every real scalar \(\zeta\), the same definition gives
\[
\|0\|_{\mathrm{coef},1}=0,\qquad
\|\zeta\|_{\mathrm{coef},1}=|\zeta|,\qquad
\|\zeta\mathcal B\|_{\mathrm{coef},1}
=|\zeta|\|\mathcal B\|_{\mathrm{coef},1},\qquad
\|x\mathcal B\|_{\mathrm{coef},1}
=\|\mathcal B\|_{\mathrm{coef},1}.
\]
Here the last identity follows by shifting the coefficient indices. Finally, expanding the product into monomials and applying the finite-sum triangle inequality yields
\[
\begin{aligned}
\mathcal B(x)\mathcal E(x)
&=\sum_{\ell,\ell'\ge0}
 \beta_\ell\gamma_{\ell'}x^{\ell+\ell'},\\
\|\mathcal B\mathcal E\|_{\mathrm{coef},1}
&\le\sum_{\ell,\ell'\ge0}
 |\beta_\ell\gamma_{\ell'}|
=\|\mathcal B\|_{\mathrm{coef},1}
 \|\mathcal E\|_{\mathrm{coef},1}.
\end{aligned}
\]
All sums here have finitely many nonzero terms.
% lean: Causalean.Mathlib.Analysis.Approximation.Chebyshev.polynomialCoeffOneNorm_mul_le

\item Define
\[
\chi:=1+\sqrt2>0,\qquad \chi^2=2\chi+1.
\]
We claim that
\[
\|\mathcal T_\ell\|_{\mathrm{coef},1}\le\chi^\ell
\qquad(\ell\in\mathbb N_0).
\]
The two initial cases are
\[
\|\mathcal T_0\|_{\mathrm{coef},1}=1=\chi^0,
\qquad
\|\mathcal T_1\|_{\mathrm{coef},1}=1\le\chi.
\]
For the induction step, the Chebyshev recurrence and the norm identities just established give
\[
\begin{aligned}
\|\mathcal T_{\ell+2}\|_{\mathrm{coef},1}
&=\|2x\mathcal T_{\ell+1}-\mathcal T_\ell\|_{\mathrm{coef},1}\\
&\le2\|\mathcal T_{\ell+1}\|_{\mathrm{coef},1}
 +\|\mathcal T_\ell\|_{\mathrm{coef},1}\\
&\le2\chi^{\ell+1}+\chi^\ell
=\chi^\ell(2\chi+1)
=\chi^{\ell+2}.
\end{aligned}
\]
This proves the claim by two-step induction.
% lean: Causalean.Mathlib.Analysis.Approximation.Chebyshev.polynomialCoeffOneNorm_chebyshev_le

\item For every \(\ell\in\mathbb N_0\),
\[
4(\ell+1)^2-1\ge3,\qquad
\frac4\pi\le\frac43,
\]
where the second inequality uses \(\pi>3\). Hence
\[
\left|
 \frac{(4/\pi)(-1)^\ell}{4(\ell+1)^2-1}
\right|
=\frac{4/\pi}{4(\ell+1)^2-1}
\le\frac12.
\]
Also, \(2\le(3/2)^2\) implies \(\sqrt2\le3/2\), and therefore
\[
\chi^2=3+2\sqrt2\le6.
\]
The product bound and the Chebyshev norm estimate now imply
\[
\begin{aligned}
\left\|
 \frac{(4/\pi)(-1)^\ell}{4(\ell+1)^2-1}
 \mathcal T_{2(\ell+1)}
\right\|_{\mathrm{coef},1}
&\le\frac12\chi^{2(\ell+1)}\\
&\le\frac12\,6^{\ell+1}
=3\cdot6^\ell
\le3\cdot8^\ell.
\end{aligned}
\]
% lean: Causalean.Mathlib.Analysis.Approximation.Chebyshev.AbsoluteValue.absChebPoly_coeffOneNorm_le

\item We next prove, by induction on \(\ell\in\mathbb N_0\), that
\[
\|q_{2\ell}\|_{\mathrm{coef},1}\le8^\ell.
\]
At \(\ell=0\),
\[
\|q_0\|_{\mathrm{coef},1}=\frac2\pi\le1=8^0.
\]
The truncations satisfy
\[
q_{2(\ell+1)}
=q_{2\ell}
+\frac{(4/\pi)(-1)^\ell}{4(\ell+1)^2-1}
 \mathcal T_{2(\ell+1)}.
\]
Thus the induction hypothesis and the preceding mode bound give
\[
\begin{aligned}
\|q_{2(\ell+1)}\|_{\mathrm{coef},1}
&\le\|q_{2\ell}\|_{\mathrm{coef},1}
+\left\|
 \frac{(4/\pi)(-1)^\ell}{4(\ell+1)^2-1}
 \mathcal T_{2(\ell+1)}
 \right\|_{\mathrm{coef},1}\\
&\le8^\ell+3\cdot8^\ell
\le8^{\ell+1}.
\end{aligned}
\]
In particular,
\[
\|q_D\|_{\mathrm{coef},1}
\le8^{\ell_{\mathrm{tr}}}
=2^{3\ell_{\mathrm{tr}}}
=2^{3D/2}.
\]
% lean: Causalean.Mathlib.Analysis.Approximation.Chebyshev.AbsoluteValue.absChebPoly_coeffOneNorm_le

\item Fix an integer \(v\) with \(0\le v\le D\). If \(a_{Dv}=0\), then
\[
|a_{Dv}|=0\le2^{3D/2}.
\]
Otherwise, its absolute value is one of the nonnegative summands in the coefficient norm. Using the stated monomial expansion of \(q_D\), we obtain
\[
|a_{Dv}|
\le\sum_{\ell=0}^{D}|a_{D\ell}|
=\|q_D\|_{\mathrm{coef},1}
\le2^{3D/2}.
\]
Together with the approximation bound, this proves both conclusions.
% lean: Causalean.Mathlib.Analysis.Approximation.Chebyshev.AbsoluteValue.absChebPoly_halfDegree_coeff_le
\end{enumerate}
\end{proof}
